\chapter{Loopback REST API}\label{rest_api}
\index{REST API}
\index{REST API!/api/version@\texttt{/api/version}}
\index{REST API!/api/load@\texttt{/api/load}}
\index{REST API!/api/model@\texttt{/api/model}}
\index{REST API!/api/randomize@\texttt{/api/randomize}}
\index{REST API!/api/train@\texttt{/api/train}}
\index{REST API!/api/train/status@\texttt{/api/train/status}}
\index{REST API!/api/train/stop@\texttt{/api/train/stop}}
\index{REST API!/api/stats@\texttt{/api/stats}}
\index{REST API!/api/regress@\texttt{/api/regress}}
\index{REST API!/api/dfa@\texttt{/api/dfa}}
\index{REST API!/api/obd@\texttt{/api/obd}}
\index{REST API!/api/cv@\texttt{/api/cv}}
\index{REST API!/api/save@\texttt{/api/save/:what}}
\index{REST API!asynchronous jobs}
\index{REST API!cancellation}
\index{REST API!busy response}
\index{REST API!action log}
\index{REST API!error response}
\index{Dataset loading!REST API}
\index{Model construction!REST API}
\index{Training!REST API}
\index{Stepwise regression!REST API}
\index{Discriminant analysis!REST API}
\index{Optimal Brain Damage!REST API}
\index{Cross-validation!REST API}
\index{Locked test!REST API}
\index{Group-disjoint partition!REST API}
\index{Covariate stratification!REST API}
\index{Sampling unit!REST declaration}
\index{Session artifacts!download API}

The graphical interface is a client of the same HTTP endpoints that
are available to scripts.  Start the service without opening a browser:

\begin{verbatim}
./build/neuron --gui --no-browser
\end{verbatim}

The engine prints a URL such as
\texttt{http://127.0.0.1:54321/}.  The port is assigned at runtime, so
clients should read it from standard output.  Requests use ordinary
form fields or multipart file uploads.  Successful JSON responses
contain \texttt{"ok":true}; failures contain
\texttt{"ok":false} and a human-readable \texttt{message}.

The server deliberately binds only to loopback and has no
authentication.  It is intended for a user or local agent on the same
computer.  It must not be exposed as a public web service.

\section{How a field's value is read}
\label{sec:rest-field-parsing}
\index{REST API!field parsing}
\index{REST API!accepted boolean tokens}
\index{REST API!omitted and empty fields}

A request field arrives as text, and the endpoint has to decide what
number, count or flag it names.  Until 2026-08-03 that decision was made
by the C conversions \texttt{atol} and \texttt{atof}, which consume as
much of the text as they understand, stop, and report nothing.  A caller
who wrote \texttt{folds=5junk} got five folds; \texttt{fraction=abc} got
zero, which is a request for a train/test split that produces no test
set and still answers \texttt{"ok":true}; and \texttt{async=true} ran
synchronously, because the only token compared was \texttt{1}.  A local
agent driving this API has no prompt to re-read and no chance to notice.

Every field is now read through the strict parsers described in
Section~\ref{sec:util-strict-parsing}, and the rules below are the whole
contract.

\subsection*{Whitespace and completeness}
Leading and trailing spaces and tabs are removed, and the rest of the
text must be consumed entirely.  \texttt{fraction=~0.3~} is accepted,
because the old conversions accepted it too; \texttt{fraction=0.3junk}
is refused.  There is no partial reading of any field.

\subsection*{Whole numbers}
Optional sign \texttt{+}, then decimal digits.  No hexadecimal, no
exponent, no decimal point.  A negative value is refused as a sign
error rather than wrapping to an enormous unsigned one, and a value
above the width of the type is refused as out of range rather than
truncating: before this change \texttt{maxiter=4294967297} trained for
one iteration and then reported that it had reached the maximum
iteration count.

\subsection*{Numbers}
Anything \texttt{strtod} accepts in the C locale, consumed entirely, and
finite.  \texttt{nan}, \texttt{inf} and their signed and mixed-case
spellings are refused for every field: an infinite gradient limit is a
stopping rule that can never fire, and an infinite variate bound is not
a bound.  Overflow to infinity is refused as out of range; underflow
yields the closest representable value, which may be zero, and the
field's own domain check then decides whether zero is allowed.

\subsection*{Flags}
Exactly \texttt{1} and \texttt{0}, case-sensitively, on every endpoint.
No other spelling is accepted anywhere, so a misspelled flag can no
longer silently mean false.  Until 2026-08-03 the five procedure flags
of \texttt{/api/cv} also accepted lowercase \texttt{true}; they no
longer do, and that is the one previously valid request this change
refuses.

\subsection*{Omitted, and present but empty}
A field that is absent leaves the value it configures alone: the
model's current setting, the engine's own default, or the documented
default for that endpoint.  A field that is present with an empty value
does the same for numbers and counts, because the page submits several
of them unconditionally --- \texttt{fraction}, \texttt{seed} and
\texttt{decay} arrive on every request whether or not the user filled
them in.

Two exceptions, both deliberate.  On \texttt{/api/train} a present but
empty \texttt{minerr}, \texttt{change}, \texttt{errwindow} or
\texttt{gradmax} \emph{disables} that stopping condition, while an
empty \texttt{printcount} is simply ignored.  And an empty
\emph{flag} is refused: there is no field on this API where a blank
checkbox means anything.

\subsection*{What a refusal looks like}
A malformed field is refused with \texttt{"ok":false} and a message
naming the field and quoting the text it was given, so that a caller can
tell which of twenty parameters was wrong:

\begin{verbatim}
folds: '5junk' is not a whole number
maxiter: '4294967297' is out of range
test_n: '-5' cannot be negative
gradmax: 'inf' is not a finite number
async: 'true' must be 1 or 0
logistic is empty
\end{verbatim}

A syntax fault and a domain fault read differently on purpose.
\texttt{strata\_bins=abc} once reported ``strata\_bins must be at least
2'', which sends the reader looking for a value problem in a value that
was never read; it now reports that the text is not a whole number,
while \texttt{strata\_bins=1} still reports the domain.

\subsection*{Refusal is atomic}
Every handler reads all of its fields before it changes anything, so a
request refused for its twentieth parameter has not applied its first.
A refused \texttt{/api/train} carrying a new learning rate leaves the
model training at the old one.

\section{A complete minimal session}

\begin{verbatim}
URL=http://127.0.0.1:54321

curl -s -X POST "$URL/api/load" \
  -d "mode=raw&path=data.txt&fraction=0.25&seed=42"

curl -s -X POST "$URL/api/model" \
  -d "type=logistic"

curl -s -X POST "$URL/api/train" \
  -d "algorithm=1&maxiter=20000&seed=42"

curl -s "$URL/api/stats"
curl -s "$URL/api/save/network" -o network.txt
curl -s "$URL/api/save/scales" -o scales.txt
\end{verbatim}

File fields accept either a path for a local script or an uploaded
multipart part for a browser.  Uploaded files are copied into the run
directory using the safe final component of their filename.

\section{Endpoint summary}

\begin{longtable}{|>{\raggedright\arraybackslash}p{0.45in}|>{\raggedright\arraybackslash}p{1.35in}|>{\raggedright\arraybackslash}p{2.35in}|}
\hline
\textbf{Verb} & \textbf{Path} & \textbf{Purpose} \\ \hline
\endfirsthead
\hline
\textbf{Verb} & \textbf{Path} & \textbf{Purpose} \\ \hline
\endhead
GET & \texttt{/api/version} & Engine package and version string. \\ \hline
POST & \texttt{/api/load} & Load, normalize, and optionally split a dataset. \\ \hline
POST & \texttt{/api/model} & Create a model or load a saved network. \\ \hline
POST & \texttt{/api/randomize} & Reset the current model's weights from a seed. \\ \hline
POST & \texttt{/api/train} & Train synchronously or start an asynchronous run. \\ \hline
GET & \texttt{/api/train/status} & Read progress and the completed result of the active or most recent long job. \\ \hline
POST & \texttt{/api/train/stop} & Request graceful cancellation at an iteration boundary. \\ \hline
GET & \texttt{/api/stats} & Return the current trained model's ROC and classification statistics. \\ \hline
POST & \texttt{/api/regress} & Run forward or reverse stepwise regression. \\ \hline
POST & \texttt{/api/dfa} & Run linear or quadratic discriminant analysis. \\ \hline
POST & \texttt{/api/obd} & Start asynchronous hidden-layer sizing. \\ \hline
POST & \texttt{/api/cv} & Start asynchronous procedure comparison and optional locked-test inference. \\ \hline
GET & \texttt{/api/save/:what} & Download a session artifact. \\ \hline
\end{longtable}

\section{Dataset loading: POST /api/load}

\begin{description}
\item[\texttt{mode}] \texttt{raw} to normalize and split one file;
\texttt{train} to load an already normalized training set and optional
matched test set.
\item[\texttt{path} or \texttt{file}] Dataset path or multipart upload.
\item[\texttt{testpath} or \texttt{testfile}] Optional matched test set
in \texttt{train} mode.
\item[\texttt{inputs}, \texttt{outputs}] Optional column-count
overrides.  Outputs default to 1; multiple outputs select BackProp and
accuracy-only reporting.
\item[\texttt{discrete}] \texttt{1} for classification or \texttt{0}
for a continuous outcome.
\item[\texttt{fraction}, \texttt{test\_n}] Test share or exact test
count.  Use one form.  A value of 0 keeps all rows in training.
\item[\texttt{val\_fraction}, \texttt{val\_n}] Optional validation
share or count for a three-way split.  The count form accompanies
\texttt{test\_n}.
\item[\texttt{seed}] Reproducible split seed.
\item[\texttt{strata}, \texttt{strata\_bins}] Comma-separated
one-based input columns for covariate stratification and the number of
quantile bins for continuous columns.
\item[\texttt{group}] Comma-separated one-based input columns whose
joint values define indivisible clusters.  Grouping takes precedence
over covariate stratification.
\item[\texttt{threshold}] Classification threshold between 0 and 1.
\item[\texttt{in\_lower}, \texttt{in\_upper}] Input normalization
bounds.
\item[\texttt{out\_lower}, \texttt{out\_upper}] Output normalization
bounds for a continuous outcome.
\item[\texttt{roc\_report}] \texttt{both} or \texttt{either}, defining
the class-availability rule for ROC reporting.
\item[\texttt{roc\_min}] Minimum observations needed for statistical
ROC reporting.
\item[\texttt{history}] \texttt{1} or \texttt{0} for dataset-operation
logging.
\end{description}

Outcome stratification is always applied to ordinary random splits.
Covariate stratification and grouping are deliberate alternatives.
A three-way split cannot currently be combined with either of them.
The response includes achieved row counts and, when appropriate, a
representativeness or zero-leakage diagnostic.

\section{Model construction: POST /api/model}

\begin{description}
\item[\texttt{type}] \texttt{logistic} or \texttt{simpleprop}.
\item[\texttt{hidden}] A comma-separated layer list such as
\texttt{5} or \texttt{8,4}.  Several layers create BackProp.
\item[\texttt{bias}] \texttt{0} creates BareProp for a one-hidden-layer
network; \texttt{1} keeps bias nodes.
\item[\texttt{errfunc}] \texttt{auto}, \texttt{xentropy}, or
\texttt{lms}.
\item[\texttt{mode}] \texttt{load} reads the network type and
architecture from a saved network.
\item[\texttt{path} or \texttt{file}] Saved network path or multipart
upload when \texttt{mode=load}.
\item[\texttt{log\_lastop}, \texttt{log\_history}] \texttt{1} or
\texttt{0} logging controls.
\end{description}

Model construction requires a loaded dataset because the input and
output dimensions are data-dependent.  A loaded network is checked
against those dimensions before it becomes current.

\section{Weight reset: POST /api/randomize}

The optional \texttt{seed} field resets the current trainable model to
a reproducible starting point.  Training itself continues from current
weights; repeated calls to \texttt{/api/train} do not randomize
implicitly.

\section{Training: POST /api/train}

\begin{description}
\item[\texttt{algorithm}] \texttt{1} for canonical gradient descent,
\texttt{2} for conjugate gradient, \texttt{3} for Shanno, \texttt{4} for
L-BFGS, or \texttt{auto} to probe canonical, conjugate gradient, and Shanno.
L-BFGS is neural-only, requires \texttt{batch\_epoch=1} and
\texttt{autostep=0}, and is not yet part of \texttt{auto}.
\item[\texttt{lbfgs\_memory}] Optional positive L-BFGS history length $m$;
accepted only with \texttt{algorithm=4}. Omission retains the current
configuration (five on a new model). Working history resets at each run.
\item[\texttt{maxiter}] Required positive iteration ceiling.
\item[\texttt{seed}] Seed used by training and automatic selection.
\item[\texttt{async}] \texttt{1} returns immediately and publishes
progress through \texttt{/api/train/status}.
\item[\texttt{eta}, \texttt{autostep}] Manual learning rate and
automatic-step toggle.
\item[\texttt{batch\_epoch}] Batch or epoch mode.  Logistic regression
requires it to remain on.
\item[\texttt{weight\_decay}, \texttt{decay}] Weight-decay toggle and
lambda.
\item[\texttt{minerr}, \texttt{change}, \texttt{errwindow},
\texttt{gradmax}] Classic stopping conditions.  Omitting a field keeps
its current value; sending an empty field disables that condition.
\item[\texttt{autostop}, \texttt{autostop\_tol},
\texttt{autostop\_window}] Plateau detector and its tolerance and
moving-window width.
\item[\texttt{logprint}, \texttt{printcount}] Iteration-table
presentation only.  These fields never change the fitted endpoint.
\end{description}

A blocking call returns the complete report, ROC series, and structured
statistics.  An asynchronous call returns acceptance immediately; the
same complete result later appears in status.  The response separates
operation success from fit validity with fields including
\texttt{converged}, \texttt{ceilingExhausted}, and
\texttt{stopReason}. It also returns the one-based \texttt{algorithm}, stable
\texttt{algorithmName}, and, for L-BFGS, \texttt{lbfgsMemory}.

\section{Long-job control}

\subsection{GET /api/train/status}

Training, stepwise regression, OBD, and cross-validation share one job
channel.  Status returns \texttt{running}, the current phase, decimated
training and held-out error series, and procedure-specific progress.
OBD adds its phase and hidden-node count.  Stepwise adds the current
pass and candidate.  Cross-validation adds stage, procedure, run,
fold, and nested-search detail when present.  When the job finishes,
\texttt{result} contains the same object its blocking form would have
returned.

\subsection{POST /api/train/stop}

Cancellation is graceful: the engine finishes the current iteration
and produces the appropriate partial audit.  Bodyless POST requests
must still carry a zero-length body:

\begin{verbatim}
curl -s -X POST -d "" "$URL/api/train/stop"
\end{verbatim}

A bare \texttt{curl -X POST} sends no Content-Length header and can
make the server wait for its read timeout before dispatching.

Only status and stop remain available while a job owns the engine.
Other engine-touching endpoints return HTTP 409 with
\texttt{"busy":true}; there is no hidden request queue.

\section{Standalone analyses}

\subsection{POST /api/dfa}

\texttt{type=linear} runs LDFA and
\texttt{type=quadratic} runs QDFA.  Optional
\texttt{log\_lastop} and \texttt{log\_history} fields control logging.
The response includes a captured report and, for a discrete
single-output dataset, ROC series and structured statistics.  DFA does
not replace the current model.

\subsection{POST /api/regress}

\begin{description}
\item[\texttt{structure}] Semicolon-separated conceptual variables;
commas join nodes and hyphens denote inclusive ranges, for example
\texttt{0;1-4;5;6,7}.
\item[\texttt{direction}] \texttt{forward} or \texttt{reverse}.
\item[\texttt{threshold}] Selection p-value threshold.
\item[\texttt{async}] \texttt{1} for status, progress, and Stop.
\end{description}

Stepwise requires a trained cross-entropy model.  Candidate subnetworks
are fitted on clones, never mutate the current model, and must converge
before entering a Wilks comparison.  The result contains the selection
path, every candidate audit, the final variables settled by completed
passes, and a \texttt{complete} flag.

\subsection{POST /api/obd}

\begin{description}
\item[\texttt{hidden\_start}, \texttt{hidden\_max}] Starting size and
safety ceiling.
\item[\texttt{iter\_budget}, \texttt{sample\_every}] Per-size ceiling
and held-out sampling cadence.
\item[\texttt{early\_stop\_tol}, \texttt{early\_stop\_patience}]
Within-size held-out reversal rule.
\item[\texttt{grow\_patience}, \texttt{prune\_tol}] Growth and pruning
rules.
\item[\texttt{autostop\_tol}, \texttt{autostop\_window}] Per-size
training plateau backstop.
\item[\texttt{algorithm}] \texttt{1}, \texttt{2}, \texttt{3}, or
\texttt{auto}.
\item[\texttt{seed}] Reproducibility seed.
\end{description}

OBD is asynchronous-only and requires a discrete, single-output
dataset with a held-out monitoring set.  Its result names that set as
\texttt{validation} or \texttt{test}, records every size trial, and
replaces the current model only after a successful search.

\subsection{POST /api/cv}

\begin{description}
\item[\texttt{folds}, \texttt{seed}] Fold count and reproducibility
seed.
\item[\texttt{logistic}, \texttt{ldfa}, \texttt{qdfa},
\texttt{neural}] Procedure-selection toggles.
\item[\texttt{neural\_obd}, \texttt{neural\_hidden}] Nested OBD or a
fixed neural size.
\item[\texttt{maxiter}, \texttt{hidden\_max},
\texttt{iter\_budget}] Ordinary-fit and nested-search ceilings.
\item[\texttt{algorithm}] Nested OBD optimizer, including
\texttt{auto}.
\item[\texttt{autostop\_tol}, \texttt{autostop\_window},
\texttt{inner\_val}] Nested-fit plateau rule and inner validation
share.
\item[\texttt{strata}, \texttt{strata\_bins}] Optional one-based
input columns for outcome-by-covariate fold balance and the number of
quantile bins for continuous columns.
\item[\texttt{group}] Optional one-based input columns whose joint
values define indivisible sampling groups.  The same key governs the
outer folds, locked holdout, and nested OBD inner validation split.
\item[\texttt{locked\_fraction}, \texttt{locked\_n}] Optional locked
test share or count.  Use one form.
\item[\texttt{primary}, \texttt{reference}] Prespecified procedure
tokens for the signed AUC contrast.
\item[\texttt{independence}] \texttt{rows} opts into ordinary paired
DeLong inference for a row-wise design.  \texttt{cluster} selects
Obuchowski clustered covariance and requires \texttt{group} plus a
locked test.  Omit the field for point estimates only.
\end{description}

Cross-validation is asynchronous-only, requires raw discrete
single-output data, and does not alter the current model.  The result
contains headline and detailed reports plus paths to Tier 3 CSV and
JSON artifacts.

Without \texttt{strata} or \texttt{group}, folds are
outcome-stratified.  \texttt{strata} balances outcome-by-covariate
cells; \texttt{group} makes whole-cluster folds for evaluation on
unseen groups.  The two cannot currently be combined.  These fields
belong to the CV request and are never inherited from
\texttt{/api/load}.  A grouped request also makes the locked holdout
and nested inner validation group-disjoint, with no row-wise fallback
when a whole-group partition is infeasible.

A locked test adds original-row predictions and a signed,
prespecified AUC contrast.  Independent-row DeLong and clustered
Obuchowski inference are explicit alternatives selected by the
declared sampling unit; neither ever falls back to the other.  The
clustered preflight requires at least two locked-test clusters carrying
each outcome class.  Tier 1 names the estimator and the number of
independent clusters.  \texttt{cv\_run.json} records the typed fold and
locked-test designs, including per-fold rows, events, groups,
imbalance, largest group, and leakage.  The two prediction CSV files
use the original raw-row identity and, for grouped designs, stable
cluster identifiers, so they join without reconstructing either key.

\section{Statistics and downloads}

\subsection{GET /api/stats}

Returns structured train and optional test statistics for the current
trained classification model.  Each set includes classification
measures, ROC estimates, confidence intervals, and curve data.  The
training endpoint also embeds these values with the captured
human-readable report.

\subsection{GET /api/save/:what}

Replace \texttt{:what} with one of:

\begin{center}
\begin{tabular}{ll}
\texttt{network} & saved model weights and architecture \\
\texttt{scales} & natural-to-normalized scaling factors \\
\texttt{train\_set}, \texttt{test\_set} & normalized datasets \\
\texttt{train\_guesses}, \texttt{test\_guesses} & model predictions \\
\texttt{dfa\_train\_guesses}, \texttt{dfa\_test\_guesses} & DFA predictions \\
\texttt{report} & captured report from the last training operation \\
\end{tabular}
\end{center}

The engine writes the artifact in the run directory and returns the
same bytes as a download.  Missing prerequisites return HTTP 400 with
a JSON explanation.

\section{Response and audit conventions}

The API returns machine-readable values together with the engine's
captured report, rather than asking clients to reconstruct statistical
meaning from a few scalars.  JSON numeric fields that cannot be
calculated are null, not invented sentinel estimates.

Every API action, including save and stop, is timestamped in
\texttt{neuron\_actions.log}.  Engine operations and self-describing
training headers are written to \texttt{neuron.log}.  Together with
the seed and saved artifacts, these logs make a local GUI or scripted
session reproducible and auditable.
